\begin{document}
\docTitle{Tailscale setup}
\phantomsection
\label{sec:tailscale-setup}
The Tailscale setup depends on whether the machine has administrator privileges.
In this MacBook configuration, a userspace \code{tailscaled} daemon runs with the following flags:
\begin{DocCode}
    --tun=userspace-networking
    --state=$HOME/.tailscale/tailscaled.state
    --socket=$HOME/.tailscale/tailscaled.sock
    --socks5-server=127.0.0.1:11055
    --outbound-http-proxy-listen=127.0.0.1:11055
\end{DocCode}
To check whether the daemon is running, use:
\begin{DocCode}
    tailscale --socket "$HOME/.tailscale/tailscaled.sock" status
\end{DocCode}
The following LaunchAgent starts the daemon automatically:
\begin{DocCode}
    ~/Library/LaunchAgents/io.tailscale.tailscaled.userspace.plist
\end{DocCode}

Commands such as \code{curl} do not use this userspace network path by default.
To route a command through it, use the local SOCKS5 proxy, for example:
\begin{DocCode}
    curl --proxy socks5h://127.0.0.1:11055 https://example.com
\end{DocCode}
\end{document}
